\subsection{Simplifying Radicals with One Variable} 
If we have a root of a power of one variable, we can simplify using the Root of a Power Property.

If the power of the variable is a multiple of the index of the root, we can proceed as follows.
\begin{example}
Simplify $\sqrt[3]{x^{18}}$.
\begin{lpic}{}{2}
\regtextleft{0.25}{-1}{Note that $18={}$};
\regtextleft{\value{cols}+0.25}{-1}{So $x^{18}={}$};
\regtextleft{0.25}{-2}{$\sqrt[3]{x^{18}}={}$};
\end{lpic}
We (do \textbar\ do not) need an absolute value because \makeblanksmall is \makeblank[0.75in].
\end{example}
\begin{example}
Simplify $\sqrt[4]{x^{20}}$.
\begin{lpic}{}{2}
\regtextleft{0.25}{-1}{Note that $20={}$};
\regtextleft{\value{cols}+0.25}{-1}{So $x^{20}={}$};
\regtextleft{0.25}{-2}{$\sqrt[4]{x^{20}}={}$};
\end{lpic}
We (do \textbar\ do not) need an absolute value because \makeblanksmall is \makeblank[0.75in].
\end{example}
If the power of the root is \emph{not} a multiple of the index, we must use division with remainder.
\begin{example}
Simplify $\sqrt[5]{x^{17}}$.
\begin{lpic}{}{2}
\regtextleft{0.25}{-1}{$17\div 5={}$};
\regtextleft{\value{cols}*2/3}{-1}{$17={}$};
\regtextleft{\value{cols}*4/3}{-1}{$x^{17}={}$};
\regtextleft{0.25}{-2}{$\sqrt[5]{x^{17}}={}$};
\end{lpic}
We (do \textbar\ do not) need an absolute value because \makeblanksmall is \makeblank[0.75in].
\end{example}
\begin{example}
Simplify $\sqrt[6]{x^{17}}$.
\begin{lpic}{}{2}
\regtextleft{0.25}{-1}{$17\div 6={}$};
\regtextleft{\value{cols}*2/3}{-1}{$17={}$};
\regtextleft{\value{cols}*4/3}{-1}{$x^{17}={}$};
\regtextleft{0.25}{-2}{$\sqrt[6]{x^{17}}={}$};
\end{lpic}
We (do \textbar\ do not) need an absolute value because \makeblanksmall is \makeblank[0.75in].
\end{example}
\begin{note}
If the power of the varaible is \makeblank[1in] than the index, we can't simplify further.
\end{note}